% !TeX encoding = UTF-8
% !TeX program = xelatex
% 中文期刊 LaTeX 模板 — 通用参考
% 基于 CTeX-org/chinesejournal 框架 (cjc.cls)
% 编译: latexmk -xelatex template.tex

\documentclass{cjc}

\usepackage{booktabs}
\usepackage{algorithm}
\usepackage{algorithmic}
\usepackage{siunitx}
\usepackage{graphicx}
\usepackage{amsmath,amssymb}

\classsetup{
  title        = {论文标题},
  title*       = {Paper Title},
  authors      = {
    author1 = {
      name         = {作者一},
      name*        = {First Author},
      affiliations = {aff1},
      biography    = {作者简介（投稿时填写）。},
      biography*   = {Author biography.},
      email        = {email@example.com},
      phone-number = {……},
    },
    author2 = {
      name         = {作者二},
      name*        = {Second Author},
      affiliations = {aff2},
      biography    = {作者简介（投稿时填写）。},
      biography*   = {Author biography.},
      email        = {email@example.com},
    },
    author3 = {
      name         = {通讯作者},
      name*        = {Corresponding Author},
      affiliations = {aff1},
      biography    = {作者简介（投稿时填写）。},
      biography*   = {Author biography.},
      email        = {email@example.com},
      corresponding = true,
    },
  },
  affiliations = {
    aff1 = {
      name  = {单位全名\ 部门全名, 城市 邮政编码, 国家},
      name* = {Department, University, City ZipCode, Country},
    },
    aff2 = {
      name  = {单位全名\ 部门全名, 城市 邮政编码, 国家},
      name* = {Department, University, City ZipCode, Country},
    },
  },
  abstract     = {中文摘要内容。},
  abstract*    = {English abstract.},
  keywords     = {关键词1, 关键词2, 关键词3, 关键词4, 关键词5},
  keywords*    = {keyword1, keyword2, keyword3, keyword4, keyword5},
  grants       = {基金资助信息（如有）。},
  % received-date = {YYYY-MM-DD},
  % revised-date  = {YYYY-MM-DD},
  % clc           = {TP391},
  % doi           = {},
}

\newcommand\dif{\mathop{}\!\mathrm{d}}

\usepackage{hyperref}

\begin{document}

\maketitle

% ============================================================
% 1  引言
% ============================================================
\section{引言}

% 研究背景与意义
% 现有工作与不足
% 本文贡献（3-4 点，每点应有证据支撑）

\section{相关工作}

% 相关工作应在主题分类下对比分析，避免"作者 [X] 提出……"的罗列模式

\subsection{子领域一}

\subsection{子领域二}

% ============================================================
% 2  方法
% ============================================================
\section{方法}

\subsection{方法总览}

% 整体架构概述 + 架构图引用

\subsection{模块一}

% 输入 → 处理 → 输出 → 设计动机

\subsection{模块二}

\subsection{损失函数}

% ============================================================
% 3  实验
% ============================================================
\section{实验}

\subsection{数据集与实验设置}

% 数据集统计、评估指标、实现细节、超参数

\subsection{与基线方法的对比}

% 主结果表 + 分析

\subsection{消融实验}

% 各模块贡献分析

\subsection{效率分析}

% 参数量、推理时间等

\subsection{定性分析}

% 可视化结果

% ============================================================
% 4  结论
% ============================================================
\section{结论}

% 总结问题、方法、关键发现
% 未来工作方向

% ============================================================
% 致谢
% ============================================================
\begin{acknowledgments}
致谢内容。
\end{acknowledgments}

% ============================================================
% 参考文献
% ============================================================
\bibliographystyle{cjc}
\bibliography{references}

% ============================================================
% 附录（可选）
% ============================================================
% \newpage
% \appendix
% \section{附录标题}
% 附录内容：详细证明、公式推导、原始数据等。

% ============================================================
% 作者简介
% ============================================================
\makebiographies

% ============================================================
% 英文背景介绍（可选）
% ============================================================
% \begin{background}
% 英文背景介绍（约 400 单词）。
% \end{background}

\end{document}
